\documentclass[12pt,a4paper,openright,twoside]{book}
\usepackage[utf8]{inputenc}
\usepackage{disi-thesis}
\usepackage{code-lstlistings}
\usepackage{notes}
\usepackage{shortcuts}
\usepackage{acronym}

\school{\unibo}
\programme{Corso di Laurea in Ingegneria e Scienze Informatiche}
\title{Progettazione di un configuratore di vendite}
\author{Cavina Mattia}
\date{\today}
\subject{Basi di dati}
\supervisor{Prof.sa Annalisa Franco}
%\cosupervisor{Dott. CoSupervisor 1}
%\morecosupervisor{Dott. CoSupervisor 2}
\session{III}
\academicyear{2025-2026}

% Definition of acronyms
\acrodef{CPQ}{configure, price, quote}
\acrodef{ERP}{enterprise resource planning}
\acrodef{CRM}{customer relationship management}
\acrodef{PLM}{product lifecycle management}
\acrodef{vm}[VM]{Virtual Machine}
\acrodef{SSIS}{SQL Server Integration Services}
\acrodef{Cepi}{Cepi S.p.a.}
\acrodef{ETO}{Engineering To Order}
\acrodef{SRP}{single-principal responsability}

\mainlinespacing{1.241} % line spacing in mainmatter, comment to default (1)

\begin{document}

\frontmatter\frontispiece

\begin{abstract}	
Max 2000 characters, strict.
\end{abstract}

\begin{dedication} % this is optional
Optional. Max a few lines.
\end{dedication}

%----------------------------------------------------------------------------------------
\tableofcontents   
\listoffigures     % (optional) comment if empty
%\lstlistoflistings % (optional) comment if empty
%----------------------------------------------------------------------------------------

\mainmatter

%----------------------------------------------------------------------------------------
\chapter{Introduzione}
\label{chap:introduction}
%----------------------------------------------------------------------------------------
Il progetto di tesi è volto alla progettazione ex-novo del programma di vendita dell'azienda \ac{Cepi}, 
azienda con una produzione di tipo \ac{ETO}.
La tipologia di produzione \ac{ETO} ha portato ad un profondo studio della situazione odierna , delle sue esigenge attuali
e future tenendo aperte le possibilità di ampliamenti.
Per far fronte a tutte le particolarità aziendali si è proceduto ad una selezione di figure chiave , per mansione o ruolo ,che potessero
dare una visione dei bisogni del proprio o con altri reparti.
La fase di analisi tramite interviste e demo dell'attaule mezzo utilizzato ha fatto emergere nuove necessità 
architteturali da affrontare come il cross-platform, l'utilizzo offline del db e la necessità di costrurire offerte in modo più agile.
L'analisi inoltre ha sottolineato come il programma non debba solo sostituire l'attuale sistema, ma debba inserirsi in un contesto
di rinnovamento più ampio coordinadosi con altri sistemi aziendali e i relativi db in ottica "\ac{SRP}"
eleggendo dei master data per ogni aspetto aziendale. Il \ac{SRP} ha plasmato notevolmente la costruzione del database ricreando, dove possibile,
le strutture dei database esistenti da cui si necessitava recuperare i dati mediante reverse Engineering.
L'intero sistema è stato pensato per appoggiarsi all'infrastruttura esistente aziendale, come SQLServer,
ma studiando delle logiche di flessibilità e ampliamento dinamico.
Si è cercato e scelto inoltre un Framework per la parte grafica che potesse
soddisfare il cross-platform essendo esso nella richieste avanzate.
Le interfacce , disegnate mediante Figma %vedere se citarlo% 
sono state pensate in ottica mobile first per rendere l'utilizzo agile da dispositivi mobili.
%Introduzione e abstract verso la fine
Le difficoltà incrontate nelle scelte preliminari erano volte ala scelta più accurata dei sistemi per dare in mano al richiedente 
un sistema che non rischiasse il debito tecnico e che potesse essere usato a lungo.
La pianificazione dei primi test con gli utenti , dei rilasci delle demo e del metodo di progettazione sono state cruciali per dettare tempi e le priorità degli item 
dovendo coprire in modo esustivo almeno il core progetto.


%Write your intro here.
%\sidenote{Add sidenotes in this way. They are named after the author of the thesis}

\paragraph{Structure of the Thesis}

\note{At the end, describe the structure of the paper}

\chapter{State of the art}

\section{Situazione As-Is}
\subsection{Situazione Infrastrutturale Aziendale}
\subsection{Attuali software utilizzati}
\subsection{Organizzazione Aziendale e Flussi di lavoro}

% --- Interviste per la raccolta requisiti ---
% \subfile{./DocumentazioneRaccolta/Analisi/Raccolta_Requisiti.tex}
\section{Raccolta dei requisiti}
\subsection{Scelta dei candidati}
\subsection{Modalità di svolgimento raccolta informazioni}
\subsection{Riassunto punti emersi per reparto}

% --- Analisi dei requisiti emersi dalle interviste ---
%\subfile{./DocumentazioneRaccolta/Analisi/Analisi_Requisiti.tex}
\section{Analisi dei requisiti}
\subsection{Requisiti funzionali}
\subsection{Requisiti non funzionali}
\subsection{Requisiti futuri should-have}

\newpage
\section{Progettazione dell'architettura}
\subsection{Preambolo sulle scelte architetturali}
\subsection{Deployment}
\subsubsection{Diagramma dei Deployment}
\subsection{Componenti}
\subsubsection{Diagramma dei Componenti}
\subsection{Flussi di lavoro}
\subsubsection{Diagrammi di sequenza}

\newpage

\section{Scelte progettuali}
\subsection{Scelta modello di sviluppo}
\subsection{Suddivisione e contenuto dei rilasci}
\subsection{Timeline di progetto e successivi focus group}
\subsection{Scelta Framework ORM}
\subsection{Scelta Framework grafico} % MAUI .NET
\subsection{Scelta del pattern di progettazione} % MVVM

\newpage
% --- Progettazione del database ---
%\subfile{./DocumentazioneRaccolta/Database/Progettazione_database.tex}
\section{Progettazione del database}
\subsection{Esportazione Vecchio Db}
\subsection{Creazione regole di denominazione}
\subsection{Utilizzo di SSIS nel processo di migrazione}

\subsection{Ristrutturazione database}
\subsubsection{Eliminazione ridondanze e normalizzazione}
\subsubsection{Creazione tabelle di lookup}
\subsubsection{Aggiunta tabelle non presenti}
\subsubsection{Creazione tabelle di Back Up}
\subsubsection{Scelta di strutture Avanzate come Triggers e Stored Procedure}

\subsection{Schema logico finale}
\subsection{Popolamento database}
\subsection{Dati fissi}
\subsection{Dati fissi da vagliare}
\subsection{Dati provenienti da altri DB}
\subsubsection{Scelta dei dati}
\subsubsection{Query di sincronizzazione e tempistiche di aggiornamento}

% --- Creazione Model e View ---
% \documentclass[12pt, a4paper]{article}

% --- Pacchetti Scientifici e Formattazione ---
\usepackage{amsmath, amssymb, amsfonts, amsthm}
\usepackage{booktabs, graphicx}
\usepackage[document]{ragged2e}
\usepackage{xcolor}

% --- Codifica e Font ---
\usepackage[utf8]{inputenc}
\usepackage[T1]{fontenc}
\usepackage[italian]{babel}
\usepackage{mathptmx}

% --- Gestione Margini e Layout ---
\usepackage[includeheadfoot]{geometry}
\geometry{
    a4paper,
    top=3cm,
    bottom=3cm,
    left=3.5cm,
    right=3.5cm,
}

% --- Gestione Interlinea e Figure ---
\usepackage{setspace}
\setstretch{1.3}
\usepackage[figuresright]{rotating}
\usepackage{float}
\usepackage[pdftex, hidelinks]{hyperref}
% Document metadata

\title{Progettazione backend}
\author{Mattia Cavina}
\date{03 febbraio 2026}

\begin{document}
\maketitle
\section{Scelta del Framework}
Dopo una analisi sui vari framwork disponibili in C\# che potessero essere adatti al nostro progetto,
abbiamo deciso di utilizzare \text{Entity Framework Core} per la gestione del database.
La connesisione al database è gestita da una classe Factory che si occupa di
creare la connesisone e restituire un oggetto \text{DbConnection} che viene utilizzato per la gestione delle chiamate al database.
\section{Pro e Contro della Scelta}
\subsection{Pro}
\begin{itemize}
    \item \textbf{Integrazione con .NET}: Entity Framework Core è progettato per integrarsi perfettamente con l'ecosistema .NET, facilitando lo sviluppo di applicazioni backend.
    \item \textbf{Supporto per LINQ}: Permette di utilizzare LINQ (Language Integrated Query) per scrivere query in modo più intuitivo e leggibile.
    \item \textbf{Migrations}: Offre un sistema di migrazioni che semplifica la gestione delle modifiche al database nel tempo.
    \item \textbf{Cross-platform}: Essendo basato su .NET Core, è compatibile con diverse piattaforme, inclusi Windows, Linux e macOS.
    \item \textbf{Mappa oggetti-relazionale}: Facilita la mappatura tra le classi C\# e le tabelle del database, riducendo la quantità di codice boilerplate necessario.
\end{itemize}
\subsection{Contro}
\begin{itemize}
    \item \textbf{Performance}: In alcuni scenari, Entity Framework Core può essere meno performante rispetto a soluzioni più leggere o a query SQL scritte manualmente.
    \item \textbf{Curva di apprendimento}: Per chi non ha familiarità con i concetti di ORM (Object-Relational Mapping), potrebbe esserci una curva di apprendimento iniziale.
    \item \textbf{Controllo limitato}: A volte, l'astrazione offerta da Entity Framework Core può limitare il controllo sulle query generate, rendendo difficile ottimizzare le prestazioni in casi complessi.
    \item \textbf{Overhead}: L'uso di un ORM può introdurre un overhead aggiuntivo rispetto a soluzioni più dirette, soprattutto in applicazioni ad alte prestazioni.
\end{itemize}
\subsection{Soluzioni ai Contro}
Per mitigare i contro di Entity Framework Core, abbiamo adottato alcune strategie:
\begin{itemize}
    \item \textbf{Ottimizzazione delle query}: Utilizziamo strumenti di profiling per identificare e ottimizzare le query lente, e in alcuni casi scriviamo query SQL personalizzate quando necessario.
    \item \textbf{Utilizzo di caching}: Implementiamo meccanismi di caching per ridurre il numero di chiamate al database e migliorare le prestazioni complessive dell'applicazione.
\end{itemize}
Queste soluzioni verranno implemenatte al bisogno durante lo sviluppo del progetto, garantendo un equilibrio tra facilità d'uso e prestazioni.
\section{Struttura}
La parte del backend verrà sviluppata in un modulo separato esponendo un'API RESTful che consente al frontend di interagire con il database.
La gestione delle query verrà incampusalta in metodi di DAO dedicati che si occuperanno di eseguire le operazioni CRUD (Create, Read, Update, Delete) sui dati.
Questo incapsulamneto permetterà l'ottimizzazione delle query e la gestione centralizzata delle operazioni sul database, facilitando la manutenzione e
l'evoluzione del codice nel tempo.
\subsection{Gestione della Sicurezza dei Dati}
La sicurezza dei dati verrà garantita  attravreso un isolamento dei degli oggetti.
Le classi di accesso ai dati creati dall ORM saranno utilizzati esclusivamente per la gestione delle query e non esposti direttamente al frontend.
Al frontend verranno esposti solo oggetti di trasferimento dati (DTO) che contengono solo le informazioni necessarie per l'interazione con l'utente,
evitando di esporre dettagli sensibili o strutture interne del database.

\section{Implementazione della mappatura}
Per mappare gli oggetti del database generati automaticamente mediante \.NET in clssi di dominio, utilizzeremo la libreria Mapster.
Mapster è una libreria di mappatura leggera e ad alte prestazioni che consente di mappare facilmente tra oggetti di diversi tipi,
come ad esempio tra entità del database e DTO (Data Transfer Objects).
Utilizzando Mapster, possiamo definire le regole di mappatura in modo semplice e flessibile,
consentendo una conversione efficiente tra le classi generate da Entity Framework Core e le classi di dominio utilizzate nel nostro backend.
Questa conversione è resa possibile mediante la classe \text(MapsterConfig), implemetazione della interfaccia \text(IRegister) di Mapster,
che definisce le regole di mappatura tra le classi generate e le classi di dominio.
Possiamo quindi configurare la profondità della mappatura per gestire correttamente le relazioni tra le entità, evitando problemi di loop infiniti o di prestazioni.

\section{Creazione delle Interfacce di Accesso ai Dati}
Si è scelto di creare dei Repository divisi per "Argomento" (ad esempio, un repository per la gestione degli utenti, uno per la gestione dei prodotti, ecc.)
che si occuperanno di incapsulare le operazioni di accesso ai dati.
Questi reposytory sono stati generati dopo un paio di classi di test per verificare la correttezza delle query e
l'efficacia della mappatura tra le classi generate da Entity Framework Core e le classi di dominio. Queste ultime sono state provate mediante unit test
per garantire la correttezza delle operazioni di accesso ai dati e la coerenza tra le classi generate e le classi di dominio.
Le interfacce generate sono:
\begin{itemize}
    \item \text{IBankRepository}: per la gestione delle operazioni relative alle banche.
    \item \text{ICustomerRepository}: per la gestione delle operazioni relative ai clienti
    \item \text{IDocumentRepository}: per la gestione delle operazioni relative ai documenti, inclusa la gestione delle righe dei documenti,
    il loro versionamento e la gestione della promozione dei documenti da offerta a ordine.
    \item \text{ILanguageRepository}: per la gestione delle operazioni relative alle lingue supportate dall'applicazione.
    \item \text{ILookupRepository}: per la gestione delle operazioni relative alle tabelle di lookup, che contengono dati di riferimento utilizzati in diverse parti dell'applicazione.
    Le operazioni all'interno saranno di sola lettura, in quanto queste tabelle vengono gestite direttamente dagli amministratori del database e non devono essere modificate dall'applicazione.
    \item \text{IMessageRepository}: per la gestione delle operazioni relative ai messaggi,ovvero notifiche e feedback che vengono inviati agli utenti in risposta a determinate azioni o eventi all'interno dell'applicazione.
    \item \text{IMountingRepository}: per la gestione delle operazioni relative ai montaggi degli impianti, inclusa la gestione degli operatori coinvolti.
    \item \text{IParagraphRepository}: per la gestione delle operazioni relative ai paragrafi, che rappresentano le sezioni o i blocchi di testo all'interno dei documenti.
    \item \text{IProductRepository}: per la gestione delle operazioni relative ai prodotti, inclusa la gestione delle loro distinte basi.
    \item \text{IPropertyRepository}: per la gestione delle operazioni relative alle proprietà, che rappresentano le caratteristiche o gli attributi dei prodotti.
    \item \text{ISellConditionRepository}: per la gestione delle operazioni relative alle condizioni di vendita, che rappresentano le regole o i criteri applicati alle offerte e agli ordini.
    \item \text{IShipperRepository}: per la gestione delle operazioni relative ai spedizionieri.
    \item \text{IStaticRepository}: per la gestione delle operazioni relative ai dati statici, che rappresentano informazioni di riferimento che non cambiano frequentemente,
    come ad esempio Nazioni, Zone Geografiche, ecc.
    \item \text{IUserRepository}: per la gestione delle operazioni relative agli utenti, inclusa la gestione dei loro ruoli e permessi all'interno dell'applicazione.
\end{itemize}

\subsection{Impletazione Repository}
I repository sopra citati sono stati sviluppati con un sistema test-driven,
creando prima i test per verificare la correttezza delle operazioni di accesso ai dati e
la coerenza tra le classi generate da Entity Framework Core e le classi di dominio, e successivamente implementando i metodi necessari per far passare i test.
Questa metodologia ha permesso di garantire la qualità del codice e la correttezza
delle operazioni di accesso ai dati fin dalle prime fasi dello sviluppo,
riducendo il rischio di bug e problemi di coerenza tra le classi generate e le classi di dominio.
Dato la continua eveluzione del progetto avere una suite di test ben definita ci ha permesso di implementare nuove funzionalità
e modificare quelle esistenti con maggiore sicurezza.
Nel corso dello sviluppo delle operazioni CRUD si è reso necessario implmentare già in questa fase le tabelle di backup
per garatire il recupero dei dati storici.

\section{Tabelle Backup}
Le tabelle di backup non sono per tutte le tabelle del database, ma solo per quelle significative come 
ad esempio la tabelle dei prodotti o dei ducumenti, che rappresentano le entità principali del nostro sistema
e che subiscono frequenti modifiche.
Esse sono state pensate con una logica ben definita:
\begin{itemize}
    \item \textbf{Struttura simile alla tabella originale}: Le tabelle di backup hanno una struttura simile a quella delle tabelle originali, con l'aggiunta di campi per tracciare la data e l'ora della modifica, l'utente che ha effettuato la modifica e un campo per indicare se la riga è stata eliminata o meno.
    \item \textbf{Gestione delle modifiche}: Ad ogni operazione di update sulla tabella originale, scatta un trigger befor update che copia la riga
    originale nella tabella di backup, mantenendo traccia delle modifiche effettuate e permettendo di recuperare lo stato precedente dei dati in caso di errori o necessità di analisi storiche.
    \item \textbf{Recupero dei dati storici}: In caso di necessità, è possibile recuperare i dati storici tramite l'id della riga ed ordinarle per "LastEditDate" per ottenere la cronologia delle modifiche effettuate su quella riga specifica.
    \item \textbf{Restore dei dati}: In caso di errori o modifiche indesiderate, è possibile utilizzare le tabelle di backup
    per ripristinare i dati allo stato precedente, garantendo la sicurezza e l'integrità dei dati nel tempo.
    Verrà comunque incremetata la data di ultima modifica per tenere traccia del momento in cui è stato effettuato il restore
    e dell'utente che ha effettuato l'operazione.
\end{itemize}
\subsection{Trigger}
I trigger come le tabelle di backup sono stati gestiti inserendoli nel codice di SSIS per garantire la loro creazione automatica durante la fase di deploy del database.
In questo modo, ogni volta che viene effettuato un deploy del database, i trigger vengono creati automaticamente, garantendo la corretta gestione delle modifiche e
il tracciamento delle operazioni sui dati fin dalle prime fasi di utilizzo del sistema.
\end{document}
% \documentclass[12pt, a4paper]{article}

% --- Pacchetti Scientifici e Formattazione ---
\usepackage{amsmath, amssymb, amsfonts, amsthm}
\usepackage{booktabs, graphicx}
\usepackage[document]{ragged2e}
\usepackage{xcolor}

% --- Codifica e Font ---
\usepackage[utf8]{inputenc}
\usepackage[T1]{fontenc}
\usepackage[italian]{babel}
\usepackage{mathptmx}

% --- Gestione Margini e Layout ---
\usepackage[includeheadfoot]{geometry}
\geometry{
    a4paper,
    top=3cm,
    bottom=3cm,
    left=3.5cm,
    right=3.5cm,
}

% --- Gestione Interlinea e Figure ---
\usepackage{setspace}
\setstretch{1.3}
\usepackage[figuresright]{rotating}
\usepackage[pdftex, hidelinks]{hyperref}

% Document metadata
\title{frontend del configuratore commerciale}
\author{Mattia Cavina}
\date{18 giugno 2026}

\begin{document}

\section{Intrduzione}
Il front end verrà realizzato mediante framework .NET MAUI nativo windows e implementato in 'C\#'.
La scelta di questo framework è stata dettata dalla necessità di avere un'applicazione integrata con i sistemi Windows,
 che possa sfruttare al meglio le funzionalità del sistema operativo e garantire una migliore esperienza utente.
La seconda motivazione è la possibilità di sviluppare un'applicazione cross-platform,
che possa essere eseguita anche su altri sistemi operativi come macOS e Linux o Android,
garantendo una maggiore flessibilità e portabilità del software.

\section{Scelte tecnologiche}
\subsection{Pattern scelto}
Per l'implementazione del frontend mediante .NET MAUI è stato scelto il pattern MVVM (Model-View-ViewModel).
Questo pattern consente di separare la logica di presentazione dalla logica di business,
facilitando la manutenzione e il test del codice.
La View rappresenta l'interfaccia utente, il ViewModel gestisce la logica di presentazione e il Model
rappresenta i dati e le regole di business.

\subsection{Framework classi View}
L'ambiente di sviluppo .NET MAUI mette a disposizione più strade di implementazione delle classi View.
La prima è quella di utilizzare le classi XAML, che consentono di definire l'interfaccia utente in modo dichiarativo tramite 
un linguaggio di markup simile a HTML.
La seconda è quella di utilizzare le classi C\# con framework ".NET Multi-platform App UI (.NET MAUI) Community Toolkit",
che consentono di definire l'interfaccia utente in modo programmatico.
La scelta di utilizzare le classi C\# con il toolkit .NET MAUI Community Toolkit
è stata dettata dalla necessità di avere un maggiore controllo sull'interfaccia utente
e sulla logica di presentazione riducendo la complessità della gestione dell'interfaccia.

\section{Definizione delle interfacce utente}
Nella realizzazione del configuratore la scelta delle interfacce utente è il passo
successivo alla realizzazione dei componenti di getsione dei dati e delle regole di business.
\subsection{Interfacce utente}
Seguendo le inteviste e il precedente sofware sono emerse le seguenti interfacce utente:
\begin{itemize}
      \item Login
      \item Menu principale
      \item Schermata ricerca documenti
      \item Schermata dettaglio documenti
      \item Schermata creazione documenti \& modifica (\ Valutare se separare in tre schermate offerta, ordine, tecnica).
      \item Schermata ricerca clienti
      \item Schermata dettaglio clienti
      \item Schermata creazione clienti \& modifica (Solo campi non importati da ERP o Fusion)
      \item Schermata ricerca elenco prodotti
      \item Schermata dettaglio prodotti \& modifica (Solo campi non importati da ERP o Fusion)
      \item Schermata creazione prodotti commerciali (Solo campi non importati da ERP o Fusion)
      \item Schermata selezione prodotti commerciali da aggiungere al documento
      \item Schermata creazione distinte prodotti commerciali (Sono i MACROGRUPPI di prodotto) \& modifica.
      \item Schermata ricerca utenti
      \item Schermata dettaglio utenti
      \item Schermata creazione utenti \& modifica
      \item Schermata gestione notifiche e feedback
      \item Schermata gestione profilo utente e modifica password
  \end{itemize}
  \subsection{Realizzazione delle prime intefacce}
  Definito l'elenco delle intefacce si è proceduti con il disegnarle almeno nella loro forma base
  madiante Figma.
  Dopo averle disegnate si è fatto un focus group con i commerciali di riferimento
  per raccogliere feedback e suggerimenti.
  \subsection{Prototipo interattivo}
      Dopo aver raccolto i feedback si è proceduto con la realizzazione di un prototipo interattivo
      che permettesse di simulare il flusso di lavoro e le interazioni tra le varie schermate.
      Il prototipo è stato realizzato direttamente in .NET MAUI e permette di navigare tra le varie schermate
      e di testare lo spostamento tra le schermate e la visualizzazione dei dati già presenti nel database.

\end{document}
\newpage
\section{Creazione Model}
\subsection{Mapping delle classi}
\subsection{Struttura del package}
\subsubsection{Classi di Test}
\subsubsection{Repository}
\subsubsection{Classi ed Interfacce esposte}

\section{Creazione View}
\subsection{Utilizzo di Figma per disegno interfacce}
\subsection{Creazione delle Factory per le View}
\subsection{Scelta delle classi di viewmodel}

\section{Modulo constroller}
\subsection{Utilizzo della Dependency Injection}
\subsubsection{Utilizzo di multi Threading per la gestione dei dati}
\subsection{Scelta gestione transazione e dei commit}
\subsection{Utilizzo della Cache in memoria per la gestione di dati fissi}

\newpage
\section{Modulo configurazione prodotto}
\subsection{Scopo del modulo}
\subsection{Focus group individuazione metodologia}
\subsection{Studio sulla guida alla configurazione}
\subsection{Query utilizzate per la ricerca}

%----------------------------------------------------------------------------------------
% BIBLIOGRAPHY
%----------------------------------------------------------------------------------------

\backmatter

\nocite{*} % Remove this as soon as you have the first citation

\bibliographystyle{alpha}
\bibliography{./Reference/references.bib}

\begin{acknowledgements} % this is optional
Optional. Max 1 page.
\end{acknowledgements}

\end{document}
